\chapter{Conclusão}
\label{cap:conclusao-seguranca}

Este manual percorreu o eBPF sob a ótica da segurança. Partimos dos fundamentos
(Capítulos~\ref{cap:introducao-seguranca} a~\ref{cap:kernel-seguranca}),
passamos pelos mecanismos que o kernel oferece para controle ativo, como LSM,
\textit{cgroup hooks} e seccomp (Capítulo~\ref{cap:mecanismos}), e pelas técnicas
de detecção e auditoria (Capítulo~\ref{cap:deteccao}). Em seguida, vimos as
ferramentas que materializam essas técnicas (Capítulo~\ref{cap:ferramentas-seguranca})
e como elas se estendem a contêineres e Kubernetes (Capítulo~\ref{cap:containers}).

O Capítulo~\ref{cap:exemplos-seguranca} reuniu essas peças em três cenários de
investigação, e o Capítulo~\ref{cap:seguranca-ebpf} lembrou que o eBPF também é
superfície de ataque: os mesmos privilégios que o tornam poderoso exigem
\texttt{CAP\_BPF} restrito, o \textit{verifier} ativo e atenção às limitações de
cada ferramenta. Detectar e bloquear com eBPF só é confiável quando o próprio
mecanismo é protegido.

Os comandos do Capítulo~\ref{cap:exemplos-seguranca} estão reunidos no script
\texttt{testes/08-seguranca-exemplos.sh}, na pasta \texttt{testes/} do repositório,
para serem executados em uma máquina de laboratório com \texttt{sudo}. Quem for
adiante pode combinar o que aprendeu aqui com o Manual~1, de observabilidade,
para enviar os eventos de segurança a um painel, ou ler a documentação do Falco e
do Tetragon para levar essas regras a um ambiente real.
